% AY2021 力学 第2問 転記（問題のみ）
\section*{第2問〔II〕}

極座標表示で $r=Ae^{\theta(t)}$ と表される軌道上を，角速度 $\dot{\theta}$ 一定で運動する
質点がある。以下の問い (1) および (2) に答えよ。ただし，$\dot{\theta}$ は $d\theta/dt$ を
表し，$A$ は正の定数，$\theta(t)$ は図に示すように $x$ 軸から反時計まわりに測った角度
とする。

\begin{enumerate}[label=(\arabic*)]
  \item $0\le\theta\le2\pi$ の範囲で，$xy$ 平面上にこの質点の運動の概形を描け。
  \item この質点の速度の大きさを求めよ。
\end{enumerate}

\begin{center}
% 図：xy平面で x軸から反時計まわりに測った角 θ(t)（動径は破線）
\begin{tikzpicture}[>=stealth, scale=1.0]
  \draw[->] (-0.4,0) -- (3.6,0) node[right] {$x$};
  \draw[->] (0,-0.4) -- (0,2.6) node[above] {$y$};
  \draw[dashed] (0,0) -- (3.1,1.55);
  \draw (1.0,0) arc (0:26.6:1.0); \node at (1.35,0.3) {$\theta(t)$};
  \node[below left] at (0,0) {$O$};
  \node at (2.0,-0.55) {図};
\end{tikzpicture}
\end{center}
